% ===========================================================
%  TITLE PAGE
% ===========================================================
\begin{titlepage}
\begin{center}

{\small\textsc{Republic of Tunisia}\\
Ministry of Higher Education and Scientific Research}

\vspace{0.3cm}
\rule{\textwidth}{1.5pt}
\vspace{0.2cm}

{\large\bfseries\color{black}
University of Carthage --- École Polytechnique de Tunisie \\[4pt]}

\vspace{1.2cm}

{\Large\bfseries\color{black} MASTER'S THESIS}\\[0.3cm]
{\normalsize Submitted in partial fulfillment of the requirements for the degree of}\\[0.3cm]
{\large\bfseries\color{black}
MASTER'S IN COMPLEX AND INTELLIGENT SYSTEMS}

\vspace{1.2cm}
\rule{\textwidth}{0.8pt}
\vspace{0.5cm}

{\LARGE\bfseries\color{black}
Learning Motor Control Strategies \\[6pt]
in Musculoskeletal Systems Using \\[6pt]
Reinforcement Learning\\[6pt]
}

\vspace{0.5cm}
\rule{\textwidth}{0.8pt}

\vspace{1.5cm}
\begin{minipage}[t]{0.48\textwidth}
\raggedright
\textbf{Presented by:}\\[6pt]
{\large M.~\textit{Mohamed Naceur Hammami}}
\end{minipage}
\hfill
\begin{minipage}[t]{0.48\textwidth}
\raggedleft
\textbf{Supervised by:}\\[6pt]
Prof.~\textit{Taysir Rezgui}
\end{minipage}

\vfill
\rule{\textwidth}{0.8pt}\\[0.3cm]
{\large\bfseries Academic Year 2025--2026}

\end{center}
\end{titlepage}

% ===========================================================
%  DEDICATION
% ===========================================================
\cleardoublepage
\thispagestyle{empty}
\vspace*{5cm}
\begin{flushright}
\itshape\large
To my parents, an inexhaustible source of love and support,\\[4pt]
whose sacrifices made this work possible.\\[8pt]
To my brother and sisters, faithful companions at every moment.\\[8pt]
To my friends and colleagues, for their constant kindness.\\[8pt]
To all those affected by neurological injury, who remind us\\
that scientific research must serve human dignity.\\[16pt]
\normalsize\bfseries\color{black} I humbly dedicate this modest work to you.
\end{flushright}

% ===========================================================
%  ACKNOWLEDGEMENTS
% ===========================================================
\frontchapter{Acknowledgements}

At the completion of this work, I would like to express my deep gratitude to all the
people who contributed, directly or indirectly, to its realisation.

First and foremost, I extend my most sincere thanks to my thesis supervisor,
Professor~Taysir Rezgui, whose availability, wise advice, and rigorous guidance were
invaluable throughout this project. Their expertise in computational biomechanics has been an intellectual model that I will strive to draw inspiration from for years to come.

I also wish to express my gratitude to the entire teaching staff of the Computer
Engineering and Cyber-Physical Systems Department at the École Polytechnique de Tunisie,
who provided me with the theoretical and methodological foundations necessary for the
successful completion of this program.

I warmly thank the members of the jury for the interest they have shown in this work
and for agreeing to evaluate it, thereby devoting a valuable part of their time to its
critical reading.

My thanks also go to the \textit{open-source} scientific community --- the developers of
\MuJoCo{}, Gymnasium, Stable-Baselines3, PyTorch, Optuna, and the Python libraries used
throughout this work --- whose efforts greatly facilitated my experiments, and without
whom reproducible reinforcement-learning research would be far more difficult.

Finally, I thank my family and loved ones for their unwavering moral support and their
patience throughout this academic journey.

\bigskip\noindent\hfill{\large\bfseries\color{black} To all of you, thank you.}

% ===========================================================
%  ABSTRACT — FRENCH (EXECUTIVE SUMMARY)
% ===========================================================
\frontchapter{Executive Summary}

Because nine muscles actuate four degrees of freedom, infinitely many force
distributions produce any given arm movement, and the movement itself can never
reveal which one the controller used. Biomechanics resolves this indeterminacy
with a minimum-effort criterion; this thesis asks whether a \textit{Deep
Reinforcement Learning} (\DRL{}) controller, trained only on a reaching reward,
arrives at the same per-muscle force distribution.

The work began from a different question --- whether a trained policy could
replace the iterative Newton--Raphson muscle solver cheaply at inference time ---
and that premise did not survive measurement. A directly solved classical
load-sharing problem runs in \SI{23.9}{\micro\second} against the network's
\SI{288.8}{\micro\second}, so no speed advantage is claimed anywhere in this
work.

A four-degree-of-freedom articulated multi-body dynamic model driven by nine
muscles was developed in \MuJoCo{} \parencite{Todorov2012}. The model specifies a
Hill-type muscle representation \parencite{Hill1938, Zajac1989}, including the contractile element with
its force-length~(FL) and force-velocity~(FV) relationships, the series elastic
element~(SEE) modelling the tendon, the parallel elastic element~(PEE), pennation-angle
dynamics, and neural activation dynamics~\parencite{Thelen2003}. Numerical integration of
the tendon-fibre equilibrium equation is handled with an iterative Newton--Raphson
method. The resulting system is wrapped as a Gymnasium simulation environment.

Four \DRL{} algorithms were implemented with Stable-Baselines3
\parencite{Raffin2021} and compared: DDPG \parencite{Lillicrap2016}, TD3
\parencite{Fujimoto2018}, SAC \parencite{Haarnoja2018}, and PPO
\parencite{Schulman2017}, alongside a random policy and a PD impedance
controller as baselines. Five configurations were trained for $10^5$ steps with
three seeds each; the best configuration was additionally trained to
$3{\times}10^5$ steps and evaluated over fifty episodes.

The central analysis compares the learned muscle forces against those prescribed
by classical static optimisation on the identical movements. The learned
controller reproduces the coordination pattern of the classical solution ---
correlation $0{,}81 \pm 0{,}04$ over five independent training runs; pattern
cosine $0{,}813$ against a shuffled null of $0{,}372$ --- while expending
$1{,}91 \pm 0{,}15$ times the minimum effort required to generate the same joint
torques, the excess concentrated in co-contraction of the elbow antagonists. A
properly tuned conventional controller solving the same load-sharing problem
explicitly attains $1{,}263$, so the learned policy uses about 51~\% more effort
than is achievable rather than 91~\% more than an unattainable ideal. Raising the
reward's effort weight closes most of that gap, to $1{,}304 \pm 0{,}025$ against
that floor.

Two secondary results follow. An error plateau at \SI{0.10}{\meter} survived
three reward redesigns and a fourfold correction to the training budget;
one-at-a-time ablation over five hyperparameters identifies the discount factor
as its sole cause, recovering $106$~\% of the default-to-tuned interval where the
other four recover nothing. Its default horizon of 100 steps spanned the whole
episode against a reach completing in 25. And configuration proved to dominate algorithm choice: with every algorithm
given a correctly matched discount factor, \SI{0.060}{\meter} separates the best
off-policy method from the second, while correcting the discount factor within a
single algorithm is worth \SI{0.165}{\meter} --- about three times as much.

The results rest on three seeds for the algorithm comparison and five for the
principal analysis. No significance testing is reported. The strict success
criterion --- one that a privileged controller satisfies only $7$~\% of the time
--- is met in $0$--$20$~\% of episodes depending on seed, a twentyfold spread
under identical settings that identifies it as noise. Limitations are set out in
full in the general conclusion.

\bigskip
\noindent\textbf{Keywords:} Musculoskeletal model, \MuJoCo{}, Hill model, Static
optimisation, Muscle redundancy, Deep Reinforcement Learning, Motor control,
PPO, SAC, TD3, DDPG, Co-contraction, Discount factor.

% ===========================================================
%  ABSTRACT — ENGLISH
% ===========================================================
\begin{otherlanguage}{english}
\frontchapter{Abstract}

Nine muscles actuate four degrees of freedom in the model studied here, so
infinitely many force distributions produce any given movement and the movement
cannot reveal which was used. Biomechanics resolves this indeterminacy by a
minimum-effort criterion. This thesis asks whether a \textit{Deep Reinforcement
Learning}~(\DRL{}) controller, trained only on a reaching reward, recovers the
same per-muscle force distribution, and under what conditions on the training
objective.

The work was begun on a different premise --- that a trained policy might replace
the iterative Newton--Raphson muscle solver cheaply at inference time. That
premise did not survive measurement: a directly solved classical load-sharing
problem takes \SI{23.9}{\micro\second} against the network's
\SI{288.8}{\micro\second}. The latency argument is withdrawn in full.

A four-degree-of-freedom multi-body articulated dynamic model driven by nine
muscles was developed using the \MuJoCo{} physics engine
\parencite{Todorov2012}. The model specifies a Hill-type muscle model \parencite{Hill1938, Zajac1989} comprising: a contractile
element~(CE) governed by force-length~(FL) and force-velocity~(FV) relationships; a
series elastic element~(SEE) representing the tendon; a parallel elastic element~(PEE);
pennation-angle dynamics; and first-order activation dynamics \parencite{Thelen2003}.
The tendon-fibre equilibrium equation is numerically resolved via Newton-Raphson
iteration at each simulation step.

Four \DRL{} algorithms were implemented and benchmarked: DDPG
\parencite{Lillicrap2016}, TD3 \parencite{Fujimoto2018}, SAC \parencite{Haarnoja2018},
and PPO \parencite{Schulman2017}, alongside two reference baselines: a random policy
and a PD impedance controller. Five configurations were trained for $10^5$ steps
with three seeds each, and the best configuration additionally to
$3{\times}10^5$ steps with fifty-episode evaluation.

The central contribution is a per-muscle load-sharing comparison between the
learned solution and classical static optimisation, posed on the identical
trajectory with the identical required joint torques and verified to machine
precision. It is a conditional comparison --- static optimisation never chooses a
trajectory of its own --- and therefore bears on force distribution, not on the
realism of the movement. The
controller reproduces the coordination structure of the classical solution
(Pearson $r = 0.81 \pm 0.04$ over five seeds; pattern cosine $0.813$ against a
shuffled null of $0.372$) at a fixed, branch-free inference cost, but expends
$1.91 \pm 0.15$ times the minimum effort needed for the same joint torques, the
excess concentrated in elbow antagonist co-contraction. Raising the reward's
effort weight reduces that ratio to $1.304 \pm 0.025$ against the $1.263$ a
properly tuned conventional controller attains, for a mild accuracy cost.
Across all fifteen trained policies the discrepancy is present, ranging from
$1.21$ to $2.30\times$, and the ordering by muscular economy is approximately the
inverse of the ordering by reaching accuracy.

Two secondary results are reported. An error plateau at \SI{0.10}{\meter}
survived three complete reward redesigns and a fourfold correction to the
effective training budget. One-at-a-time ablation over five hyperparameters,
three seeds each, identifies the \emph{discount factor} as its sole cause: it
recovers $106$~\% of the default-to-tuned interval while the target entropy,
learning rate, target-network rate and batch size each recover nothing. Its
default value implies a 100-step horizon --- the entire episode --- against a
reach that completes in 25 steps. Separately, configuration was found to dominate algorithm choice: with every
algorithm given a correctly matched discount factor, \SI{0.060}{\meter} separates
the best off-policy method from the second, against \SI{0.165}{\meter} recovered
by correcting that one setting within a single algorithm.

These results rest on three seeds for the algorithm comparison and the
ablations, and five for the principal analysis and the effort study; no
significance testing is reported. Under the strict
success criterion --- which a privileged controller with full state access
satisfies only $7$~\% of the time --- the best policy scores $0$--$4$~\%.
Limitations are stated in full in the general conclusion.

\bigskip
\noindent\textbf{Keywords:} Musculoskeletal model, \MuJoCo{}, Hill muscle model,
Static optimisation, Muscle redundancy, Newton--Raphson solver, Deep
Reinforcement Learning, Motor control, PPO, SAC, TD3, DDPG,
Co-contraction, Discount factor.
\end{otherlanguage}

% ===========================================================
%  LIST OF ABBREVIATIONS
% ===========================================================
\frontchapter{List of Abbreviations}

{%
\setstretch{1.0}%
\setlength{\tabcolsep}{5pt}%
\renewcommand{\arraystretch}{0.92}%
\normalsize
\begin{tabularx}{\textwidth}{@{} L{2.8cm} X @{} }
\toprule
\rowcolor{lightblue}
\textbf{Abbrev.} & \textbf{Meaning} \\
\midrule
\rowcolor{lightblue!40}
BCa      & \textit{Bias-Corrected and accelerated (bootstrap)} \\
CCI      & Co-contraction index (Falconer \& Winter) \\
\rowcolor{lightblue!40}
CE       & \textit{Contractile Element} (Hill model) \\
CNS      & Central nervous system \\
\rowcolor{lightblue!40}
DOF      & Degrees of freedom \\
DDPG     & \textit{Deep Deterministic Policy Gradient} \\
\rowcolor{lightblue!40}
DR       & \textit{Domain Randomization} \\
DRL      & \textit{Deep Reinforcement Learning} \\
\rowcolor{lightblue!40}
EMA      & \textit{Exponential Moving Average} \\
EMG      & Electromyography \\
\rowcolor{lightblue!40}
FL       & Force-length (muscle relationship) \\
FV       & Force-velocity (muscle relationship) \\
\rowcolor{lightblue!40}
GAE      & \textit{Generalised Advantage Estimation} \\
GPU      & \textit{Graphics Processing Unit} \\
\rowcolor{lightblue!40}
MDP      & \textit{Markov Decision Process} \\
\rowcolor{lightblue!40}
MJCF     & \textit{MuJoCo XML Format} \\
MLP      & \textit{Multi-Layer Perceptron} \\
\rowcolor{lightblue!40}
\MuJoCo  & \textit{Multi-Joint dynamics with Contact} \\
NMF      & \textit{Non-negative Matrix Factorisation} \\
\rowcolor{lightblue!40}
NRMSE    & \textit{Normalised Root Mean Squared Error} \\
OFC      & \textit{Optimal Feedback Control} \\
\rowcolor{lightblue!40}
SLSQP    & \textit{Sequential Least-Squares Quadratic Programming} \\
VAF      & Variance accounted for (synergy decomposition) \\
\rowcolor{lightblue!40}
OpenSim  & \textit{Open-Source Musculoskeletal Simulator} \\
\rowcolor{lightblue!40}
OU       & Ornstein-Uhlenbeck (stochastic process) \\
PD       & \textit{Proportional-Derivative} (controller) \\
\rowcolor{lightblue!40}
PEE      & \textit{Parallel Elastic Element} (Hill model) \\
PPO      & \textit{Proximal Policy Optimisation} \\
\rowcolor{lightblue!40}
ReLU     & \textit{Rectified Linear Unit} \\
\RL{}    & \textit{Reinforcement Learning} \\
\rowcolor{lightblue!40}
SAC      & \textit{Soft Actor-Critic} \\
\rowcolor{lightblue!40}
SEE      & \textit{Series Elastic Element} --- the tendon (Hill model) \\
\rowcolor{lightblue!40}
TD3      & \textit{Twin Delayed Deep Deterministic Policy Gradient} \\
TPE      & \textit{Tree-structured Parzen Estimator} (sampler Optuna) \\
\bottomrule
\end{tabularx}
}

% ===========================================================
%  LIST OF SYMBOLS
% ===========================================================
\frontchapter{List of Symbols}

{%
\setstretch{1.0}%
\setlength{\tabcolsep}{4pt}%
\renewcommand{\arraystretch}{0.90}%
\footnotesize
\begin{tabularx}{\textwidth}{@{} L{2.15cm} L{2.0cm} X @{} }
\toprule
\rowcolor{lightblue}
\textbf{Symbol} & \textbf{Unit} & \textbf{Meaning} \\
\midrule
\multicolumn{3}{@{}l}{\textit{\color{navyblue}\textbf{Hill Muscle Model}}}\\
\midrule
$a(t)$                 & ---     & Muscle activation, $a \in [0,1]$ \\
\rowcolor{lightblue!40}
$u(t)$                 & ---     & Neural excitation, $u \in [0,1]$ \\
$\tau_\text{act}, \tau_\text{deact}$ & s & Activation/deactivation time constants \\
\rowcolor{lightblue!40}
$l^M$, $l^M_\text{opt}$ ($\equiv l_0^M$) & m & Fibre length, optimal length \\
$\lnorm = l^M/l_0^M$   & ---     & Normalised fibre length \\
\rowcolor{lightblue!40}
$l^T$, $l^T_\text{slack}$ & m    & Tendon length, slack length \\
$l^\text{MT}$          & m       & Total musculotendon-unit length \\
\rowcolor{lightblue!40}
$\alpha$, $\alpha_0$   & rad     & Current pennation angle / value at $l_0^M$ \\
$v_\text{max}$         & $l_0^M$/s & Maximum shortening velocity \\
\rowcolor{lightblue!40}
$\vnorm$               & ---     & Normalised contraction velocity \\
$f_\text{FL}, f_\text{FV}$ & --- & Force-length / force-velocity factors \\
\rowcolor{lightblue!40}
$F_0^M$                & N       & Maximum isometric force \\
$F_\text{SEE}, F_\text{PEE}$ & N & Tendon (SEE) and passive (PEE) forces \\
\rowcolor{lightblue!40}
$F^\text{muscle}$      & N       & Total muscle force applied to the skeleton \\
$k_T, k_\text{PE}$     & ---     & Dimensionless stiffness parameters (tendon, PEE) \\
\rowcolor{lightblue!40}
$\gamma_\text{FL}$     & ---     & Width parameter of the FL curve \\
\midrule
\multicolumn{3}{@{}l}{\textit{\color{navyblue}\textbf{Multi-Body Dynamics}}}\\
\midrule
$q, \dot{q}$           & rad, rad/s & Generalised joint positions / velocities \\
\rowcolor{lightblue!40}
$p_\text{ee}$          & m       & End-effector position (hand) \\
$p_\text{target}$      & m       & Target position \\
\rowcolor{lightblue!40}
$\Delta p, \|\Delta p\|$ & m     & Error vector and distance to target \\
$J(q)$                 & m/rad   & Geometric Jacobian of the end effector \\
\midrule
\multicolumn{3}{@{}l}{\textit{\color{navyblue}\textbf{Reinforcement Learning}}}\\
\midrule
$\mathcal{S}, \mathcal{A}$ & ---   & State / action space \\
\rowcolor{lightblue!40}
$\pi_\theta(a|s)$      & ---     & Policy parameterised by $\theta$ \\
$Q^\pi, V^\pi, A^\pi$  & ---     & Action-value, state-value, and advantage functions \\
\rowcolor{lightblue!40}
$\gamma$               & ---     & Discount factor ($\in [0,1)$) \\
$r(s,a)$               & ---     & Reward signal \\
\rowcolor{lightblue!40}
$\alpha$ (SAC)         & ---     & Entropy coefficient (temperature) \\
$\mathcal{H}(\pi)$     & nats    & Policy entropy \\
\rowcolor{lightblue!40}
$\tau$ (Polyak)        & ---     & Target-network update coefficient \\
$\mathcal{B}$          & ---     & Replay memory (\textit{replay buffer}) \\
\midrule
\multicolumn{3}{@{}l}{\textit{\color{navyblue}\textbf{Statistics and Evaluation}}}\\
\midrule
$N_\text{seeds}$       & ---     & Number of random seeds ($= 10$) \\
\rowcolor{lightblue!40}
$W$                    & ---     & Wilcoxon test statistic \\
$r_\text{rb}$          & ---     & Rank-biserial correlation (effect size) \\
\rowcolor{lightblue!40}
$\text{CI}_{95}$       & ---     & 95~\% confidence interval (BCa) \\
\bottomrule
\end{tabularx}
}

% ===========================================================
%  TABLE OF CONTENTS / FIGURES / TABLES
% ===========================================================
\clearpage
\tableofcontents

\clearpage
\listoffigures

\clearpage
\listoftables
